%! Rational Functions and Holes — Unit 3, Lesson 3.4 — Group Activity
\documentclass[10pt]{article}
\usepackage{precalculus-article}
\usepackage{precalculus-boxes}
\newcommand{\TallMath}[1]{$\displaystyle #1\rule[-1.4em]{0pt}{3.2em}$}

\begin{document}

\pageheader{Unit 3, Lesson 3.4}{Group Activity}
\namepartnerperiod

\vspace{0.1in}

\begin{scenariobox}[Missing Data Points]{sky}
\small
A scientist's measurement device records four data streams, each modeled by a rational function.
At certain input values, the device has a gap in the data. Your job is to analyze each function,
determine whether each gap is a \textbf{hole} or a \textbf{vertical asymptote}, and --- if it
is a hole --- find the exact coordinates.

\medskip
The four functions are:
\[
  f_1(x) = \frac{x^2-9}{x-3}
  \qquad
  f_2(x) = \frac{x^2-1}{(x-1)(x+4)}
  \qquad
  f_3(x) = \frac{(x-2)^2}{x^2-4}
  \qquad
  f_4(x) = \frac{x^2+2x}{x}
\]
\end{scenariobox}

\vspace{0.1in}

\begin{tcolorbox}[colback=white, colframe=black!40,
    title=\textbf{Tier R --- Remediate},
    fonttitle=\bfseries, arc=2mm, left=3mm, right=3mm, top=2mm, bottom=2mm]

The numerator and denominator of each function are already written in a helpful form below.
For each function, circle the shared factor, then decide: is $x = c$ a \textbf{hole} or a \textbf{VA}?

\renewcommand{\arraystretch}{2.0}
\begin{tabularx}{\linewidth}{@{} c >{\centering}p{3.8cm} >{\centering}p{3.8cm} c X @{}}
\textbf{Fn} & \textbf{Numerator (factored)} & \textbf{Denominator (factored)} & \textbf{Gap at} & \textbf{Hole or VA?} \\
\midrule
$f_1$ & $(x-3)(x+3)$ & $(x-3)$ & $x=3$ & \blank{2.5cm} \\
$f_2$ & $(x-1)(x+1)$ & $(x-1)(x+4)$ & $x=1$ & \blank{2.5cm} \\
      &              &              & $x=-4$ & \blank{2.5cm} \\
$f_3$ & $(x-2)^2$ & $(x-2)(x+2)$ & $x=2$ & \blank{2.5cm} \\
      &           &              & $x=-2$ & \blank{2.5cm} \\
$f_4$ & $x(x+2)$ & $x$ & $x=0$ & \blank{2.5cm} \\
\end{tabularx}
\end{tcolorbox}

\vspace{0.1in}

\begin{tcolorbox}[colback=white, colframe=black!40,
    title=\textbf{Tier A --- Apply},
    fonttitle=\bfseries, arc=2mm, left=3mm, right=3mm, top=2mm, bottom=2mm]

\begin{enumerate}[itemsep=10pt]
  \item Factor the numerator and denominator of each function yourself. Show your work.
        Then identify all shared zeros and apply the multiplicity rule to classify each as
        a hole or VA.

  \item For each hole you identified, find the coordinates $(c, L)$. Show the cancellation
        and substitution steps.

        \medskip
        \renewcommand{\arraystretch}{1.8}
        \begin{tabular}{|c|c|c|c|}
        \hline
        \textbf{Function} & \textbf{Shared zero} $c$ & \textbf{Simplified form} & \textbf{Hole at $(c, L)$} \\
        \hline
        $f_1$ & \blank{2cm} & \blank{3cm} & \blank{2.5cm} \\
        \hline
        $f_2$ & \blank{2cm} & \blank{3cm} & \blank{2.5cm} \\
        \hline
        $f_3$ & \blank{2cm} & \blank{3cm} & \blank{2.5cm} \\
        \hline
        $f_4$ & \blank{2cm} & \blank{3cm} & \blank{2.5cm} \\
        \hline
        \end{tabular}
\end{enumerate}
\end{tcolorbox}

\vspace{0.1in}

\begin{tcolorbox}[colback=white, colframe=black!40,
    title=\textbf{Tier E --- Extend},
    fonttitle=\bfseries, arc=2mm, left=3mm, right=3mm, top=2mm, bottom=2mm]

\begin{enumerate}[itemsep=10pt]
  \item Write the limit notation for each hole from Tier A.
        For example: $\lim_{x \to c} f_i(x) = L$.

        \medskip
        $\displaystyle\lim_{x \to \blank{0.8cm}} f_1(x) = $ \blank{1.5cm} \qquad
        $\displaystyle\lim_{x \to \blank{0.8cm}} f_2(x) = $ \blank{1.5cm} \qquad
        $\displaystyle\lim_{x \to \blank{0.8cm}} f_3(x) = $ \blank{1.5cm} \qquad
        $\displaystyle\lim_{x \to \blank{0.8cm}} f_4(x) = $ \blank{1.5cm}

  \item Choose one hole from above. Explain in a complete sentence \textit{why} the
        $y$-coordinate of the hole is that specific value, referencing what you did to find it.

        \vspace{0.06in}
        \writelines{2}

  \item Construct a rational function $g(x)$ that has \textbf{both} a hole at $x = 1$
        \textbf{and} a vertical asymptote at $x = -2$. Write your answer in factored form
        and in simplified form.

        \vspace{0.06in}
        Factored: $g(x) = $ \blank{6cm}

        \vspace{0.06in}
        Simplified (hole removed): $g(x) = $ \blank{5cm} \quad with hole at $\Big($\blank{0.8cm}$,$ \blank{1.5cm}$\Big)$
\end{enumerate}
\end{tcolorbox}

\end{document}
